\documentclass{article}

\usepackage{natbib}
\usepackage{booktabs}
\usepackage{tabularx}
\usepackage{hyperref}

\hypersetup{
    colorlinks=true,       % false: boxed links; true: colored links
    linkcolor=red,          % color of internal links (change box color with linkbordercolor)
    citecolor=green,        % color of links to bibliography
    filecolor=magenta,      % color of file links
    urlcolor=cyan           % color of external links
}

\title{Hazard Analysis\\\progname}

\author{\authname}

\date{}

\input{../Comments}
\input{../Common}

\begin{document}

\maketitle
\thispagestyle{empty}

~\newpage

\pagenumbering{roman}

\begin{table}[hp]
\caption{Revision History} \label{TblRevisionHistory}
\begin{tabularx}{\textwidth}{llX}
\toprule
\textbf{Date} & \textbf{Developer(s)} & \textbf{Change}\\
\midrule
Date1 & Name(s) & Description of changes\\
Date2 & Name(s) & Description of changes\\
... & ... & ...\\
\bottomrule
\end{tabularx}
\end{table}

~\newpage

\tableofcontents

~\newpage

\pagenumbering{arabic}

\wss{You are free to modify this template.}

\section{Introduction}

The team's definition of a hazard is an “unsafe” condition that may lead to an unintended event that causes an undesirable outcome \citep{IppolitoWallace1995}. In relation to software development, hazards can be understood as the underlying design decisions that result in potential harm. These risks become apparent when the hazards are not addressed prior to development. These hazards can include technical failures, privacy/security, usability and accessibility risks. 

The project's goal is to develop a tool that provides guidance for a physiotherapist's prescribed home exercises. The user's form is analyzed and prompts for adjustments are given. While this can increase the rate and efficacy of recovery, it introduces several hazards. \href{https://github.com/M9Huynh/technically-functional/blob/main/docs/SRS/SRS.pdf}{See the Section 3 in the Problem Statement document for more on project goals}. The purpose of this document is to identify those hazards, outline assumptions and propose solutions to limit their presence and their frequency.


\section{Scope and Purpose of Hazard Analysis}

This document outlines the hazards that may occur while operating the tool. The key areas for scope include:


\begin{itemize}
    \item Technical Hazards: Inability to process user footage, incorrect identification of exercise from different planes, or inaccurate data comparisons, which can result in performance issues and impact the functionality of the tool. Furthermore, these hazards can result in misleading feedback, which could contribute to improper and unsafe exercise techniques.
    \item Usability Hazards: Usability ensures that the tool has the potential for  successfully being adopted. All of the inputs to the system, such as the recording of the exercise, are generated from the user. A user-friendly interface guarantees ease of use, which encourages utilization of this tool. On the other hand, incorrect, on-screen prompts  may cause confusion  while the user performs the exercise. This may reduce usability and lead to ambiguous guidelines.
    \item Privacy/Security Hazards: Unnecessary data collection and storage increases the potential of unauthorized access and misuse of personal information. Such situations, such as data breaches, may lead to catastrophic outcomes, which will result in lower adoption rate.
\end{itemize}


The purpose of this document is to identify risks and their associated impact to determine ways to mitigate them and identify any future occurrences 


\section{System Boundaries and Components}

To divide the system boundaries we outline the key functions of the application.

\begin{itemize}
    \item Authentication module
    \begin{itemize}
        \item The user will log in to the application to access prior user data.
    \end{itemize}
    \item Video streaming module
    \begin{itemize}
        \item The user will video stream themselves performing a physio exercise pertaining to the knee joint.
    \end{itemize}
    \item Analysis module
    \begin{itemize}
        \item The input of the user performing an exercise will be assessed against the requirements and guidelines defined in the backend.
    \end{itemize}
    \item Reporting module
    \begin{itemize}
        \item The system will generate analytics from the analysis module which packages the information to be viewed by the user or physiotherapist.
        \end{itemize}
\end{itemize}


In addition we will consider the following within the hazard analysis: 

\begin{itemize}
    \item A device with a camera
    \begin{itemize}
        \item To allow users to record themselves during the motion and to receive feedback on the exercise performed.
    \end{itemize}
    \item A database where user data will be stored
\end{itemize}


These system boundaries outline what is covered within this hazard analysis.


\section{Critical Assumptions}

\wss{These assumptions that are made about the software or system.  You should
minimize the number of assumptions that remove potential hazards.  For instance,
you could assume a part will never fail, but it is generally better to include
this potential failure mode.}

\section{Failure Mode and Effect Analysis}

\wss{Include your FMEA table here. This is the most important part of this document.}
\wss{The safety requirements in the table do not have to have the prefix SR.
The most important thing is to show traceability to your SRS. You might trace to
requirements you have already written, or you might need to add new
requirements.}
\wss{If no safety requirement can be devised, other mitigation strategies can be
entered in the table, including strategies involving providing additional
documentation, and/or test cases.}

\section{Safety and Security Requirements}

\wss{Newly discovered requirements.  These should also be added to the SRS.  (A
rationale design process how and why to fake it.)}

\section{Roadmap}

\wss{Which safety requirements will be implemented as part of the capstone timeline?
Which requirements will be implemented in the future?}

\newpage{}

\section*{Appendix --- Reflection}

\wss{Not required for CAS 741}

\input{../Reflection.tex}

\begin{enumerate}
    \item What went well while writing this deliverable? 
    \item What pain points did you experience during this deliverable, and how
    did you resolve them?
    \item Which of your listed risks had your team thought of before this
    deliverable, and which did you think of while doing this deliverable? For
    the latter ones (ones you thought of while doing the Hazard Analysis), how
    did they come about?
    \item Other than the risk of physical harm (some projects may not have any
    appreciable risks of this form), list at least 2 other types of risk in
    software products. Why are they important to consider?
\end{enumerate}

\bibliographystyle {plainnat}
\bibliography{../../refs/References}

\end{document}